\chapter{Central Venous Port Placement}

\section*{Introduction \& Clinical Indications}
A totally implanted venous port provides repeated, reliable access for chemotherapy, prolonged infusions, transfusions, or selected patients requiring intermittent intravenous treatment. The reservoir sits beneath the skin and connects to a catheter terminating in a central vein. Because no external tubing remains between uses, ports can be useful when access is needed for months rather than days.

The decision should balance treatment duration, infusion characteristics, infection risk, venous anatomy, and patient preference. A port is not automatically appropriate for every difficult intravenous-access problem; short-term peripheral or midline access may be safer for compatible therapies. Vesicant drugs and infusates with extreme pH or osmolarity require a device and administration protocol designed for central delivery.

\section*{Relevant Surgical Anatomy}
Common approaches use the internal jugular or subclavian region, with the catheter traveling through the brachiocephalic vein and superior vena cava. The right side often gives a straighter route. The carotid artery, pleura, thoracic duct, and nearby nerves are at risk during access. The port pocket is usually made on the upper chest where it can be palpated and protected from pressure or clothing.

\section*{Preoperative Workup \& Preparation}
Review indication, expected duration, allergies, anticoagulants, platelet count, infection, prior lines, and venous thrombosis. Ultrasound guidance and fluoroscopy or another reliable method of tip confirmation are used according to local practice. Treat active bacteremia before elective implantation when possible. Discuss bleeding, infection, thrombosis, catheter fracture or migration, pneumothorax, malposition, and the need for removal if the device fails.

\section*{Surgical Technique \& Procedure}
After sterile preparation and local anesthesia with or without sedation, the operator punctures the target vein under ultrasound guidance and advances a guidewire. A subcutaneous pocket is created, the catheter is tunneled to the venous entry site, and its length is selected so the tip lies in an appropriate central position. The catheter is connected to the reservoir, flushed, and secured in the pocket. The incisions are closed and a sterile dressing applied.

The port is accessed only with a non-coring needle. Patency, blood return, and radiographic position are checked before use when required. If resistance or pain occurs, the port should not be forcibly flushed; mechanical obstruction, malposition, pinch-off, thrombosis, or extravasation must be considered.

\section*{Complications \& Risks}
Immediate risks include arterial puncture, hematoma, pneumothorax, air embolism, arrhythmia, and malposition. Delayed risks include pocket infection, bloodstream infection, thrombosis, fibrin sheath, catheter fracture, port rotation, skin erosion, and extravasation. New arm or neck swelling, fever, drainage, chest symptoms, loss of blood return, or pain during infusion warrants urgent review.

\section*{Postoperative Care \& Recovery}
Keep the incision clean and dry as directed, avoid heavy pressure over the pocket during healing, and provide written access and emergency instructions. Staff should use aseptic technique, disinfect access points, and follow institutional flushing and locking policies. Remove the port when it is no longer needed or when infection or another complication makes retention unsafe.

\section*{Outcomes \& Prognosis}
Ports can provide durable access with low day-to-day burden when implanted and maintained correctly. Outcomes depend on patient selection, sterile access, prompt recognition of infection or thrombosis, and coordination between oncology, infusion, surgery, and radiology teams.

\section*{References}
\begin{enumerate}
\item CDC. Guidelines for the Prevention of Intravascular Catheter-Related Infections.
\item Society of Interventional Radiology. Guidelines for Central Venous Access.
\item National Institute for Health and Care Excellence. Guidance on preventing healthcare-associated infections.
\end{enumerate}
